problem:  
Let \( M^{n} \) be a compact, simply connected Riemannian manifold with **almost positive sectional curvature**—that is, the subset of points \( p \in M \) for which every tangent 2-plane \( \sigma \subset T_pM \) satisfies \( K_p(\sigma) > 0 \) is open and dense in \( M \).  
Assume \( M \) admits an isometric **Lie group action** \( G \curvearrowright M \) with cohomogeneity one, and with principal orbits diffeomorphic to \( G/H \). Suppose further that \( M \) is **stably parallelizable** (i.e. \( TM \oplus \varepsilon^k \cong \varepsilon^{n+k} \) for some \( k \)).  

**Question:**  
Does there exist any stably parallelizable, compact, simply connected **cohomogeneity-one** manifold admitting a metric of almost positive sectional curvature?  
Equivalently, does the combination of almost positive curvature and cohomogeneity-one symmetry necessarily force nontrivial Pontryagin classes (or nontrivial stable tangential data)?  

Why is it a "good" problem:  
This problem corrects and refines the previous formulation by introducing a powerful geometric constraint—**cohomogeneity-one symmetry**—which makes the question both sharper and tractable. Cohomogeneity-one manifolds form a richly structured and well-classified class where curvature properties, topology, and group actions can be analyzed concretely via the **group diagram** \( H \subset K_\pm \subset G \).  

By restricting to such manifolds, we can meaningfully ask how the topological condition of **stable parallelizability** interacts with the geometric condition of **almost positive curvature**—a relaxation of positive curvature that currently has only a handful of known examples (mostly Eschenburg and Bazaĭkin spaces, none of which are stably parallelizable).  This focuses the challenge on whether the open–dense positivity of curvature already enforces nontrivial topology through Pontryagin class obstructions, or whether genuine counterexamples can exist within the controlled setting of cohomogeneity-one actions.  

The problem thereby connects three central themes:
- **Curvature:** the borderline between positive and almost positive sectional curvature.  
- **Topology:** the stability of the tangent bundle via Pontryagin classes and parallelizability.  
- **Symmetry:** cohomogeneity-one classification, which provides explicit topological models.  

Its statement is mathematically simple yet conceptually deep. It demands an understanding of how curvature-induced geometric rigidity coexists (or fails to) with topological triviality under strong symmetry. A positive resolution would extend known rigidity phenomena to the almost-positive realm; a counterexample would uncover new freedom mechanisms. The problem elegantly intertwines Riemannian geometry, transformation group theory, and characteristic class topology, making it a genuinely *good* research problem: accessible in statement, but profound in consequence.